% M10, 11.09.2026: Leserführung, Kapiteltrennung und belegte Reichweite.
% Resume: runs/fullworkflow/20260805_144056_dc711f; Ledger c939cd.
% Pause d2860ea5bd56415698b6d3d951acef32: Checkpoint-Index Position 13, stepNumber 12.
% Ereignisse STAGE_LEDGER_DONE -> PIPELINE_PAUSED -> PIPELINE_RESUME -> RUN_DONE; keine PIDs.
% Core-Rücksetzung: iteration-notes-ledger-kapsel.md, 05.08.; narrativer Beleg.
% Smoke tools/smoke-hitl.sh: je ein Resume in decision / pipeline-full, 14 Prüfungen.
% F2/F5b aus 7.5 unverändert übernommen; keine neuen Versuche oder Messergebnisse.

\section{Wiederaufnahme und agentische Bedienung}
\label{sec:eval-resume}

Die Anforderungen A9 (Durabilität) und A10 (kontrollierte agentische
Bedienung) werden durch eine Wiederaufnahme-Demonstration, technische
Prüfungen und zwei dokumentierte Steward-Aufgaben adressiert. Diese
Belege untersuchen unterschiedliche Ausschnitte des Architekturanspruchs.

\textbf{Wiederaufnahme.} Im dokumentierten Pipeline-Lauf war die
Ledger-Stufe bereits abgeschlossen, als der Lauf am Adjudikations-Gate
pausierte. Bis zu dieser Pause waren 13 Checkpoints des Elternlaufs
gespeichert. Die Ereigniskette zeigt die Wiederaufnahme desselben
Checkpoints, die Weiterverwendung der Zuordnung zum abgeschlossenen
Ledger-Teillauf und den anschließenden regulären Abschluss. Eine erneute
Ledger-Ausführung ist dabei nicht verzeichnet. Belegt ist damit dieser
Fortsetzungsfall, nicht die Wiederaufnahme innerhalb einer noch laufenden
Ledger-Verarbeitung. Die Ereignisse weisen keine Prozessidentitäten aus;
der gesonderte technische Nachweis der Prozessgrenze folgt unten.
Außenwirkungen waren in der Demonstration deaktiviert. Die
Entwicklungsnotiz berichtet die anschließende bytegleiche
Wiederherstellung des zuvor gesicherten Core. Diese Rücksetzung ist
kein Vergleich mit dem Ergebnis einer ununterbrochenen
Referenzausführung; ein solcher Vergleich wurde nicht durchgeführt.

Der deterministische \emph{Prozessgrenzen-Smoke-Test} umfasst
14 Prüfungen. Er startet unter anderem einen Entscheidungsworkflow
und die Gesamtkette über den Delta-Eingang bis zu einer Gate-Pause
und setzt beide jeweils einmal durch einen neuen CLI-Prozess fort.
Die Rücksetzung des Testzustands ist eine davon getrennte
Prüfhandlung, keine Leistung der Workflow-Checkpoints.
Der Test sichert die bezeichnete Mechanik ab; er prüft weder die
fachliche Ergebnisqualität noch die vollständige Transkriptverarbeitung.
Stände und Belegzuordnung stehen in Anhang~\ref{anh:evaluation}.

\textbf{Agentische Bedienung (Steward).} Zwei bereits in
Abschnitt~\ref{sec:eval-fortschreibung} belegte Aufgaben zeigen den
Steward über den regulären Gates und zählen nicht als zusätzliche
unabhängige Versuche: In Fall F2 wurde ein frei diktierter
Präzisierungswunsch als Eingangs-Delta gefasst, am Ingest-Gate
vorgelegt und nach Freigabe als Fortschreibung wirksam. In Fall F5b
führte ein Steward-Werkzeug externe GitHub-Kommentare in denselben
Gate-Weg. In beiden Aufgaben wählte der Agent das vorgesehene
Werkzeug; die betrachteten fachlichen Core-Änderungen erfolgten nach
Freigabe über den regulären Anwendungspfad. Eine Usability- oder
Produktivitätsstudie der Bedienschicht fand nicht statt.
